\item Un haz de luz monocromático es absorbido por un conjunto de átomos de hidrógeno en estado fundamental, de manera que se observan exactamente seis longitudes de onda diferentes en el espectro de emisión cuando el hidrógeno se relaja de nuevo al estado fundamental.
\begin{enumerate}
    \item ¿Cuál es la longitud de onda del haz incidente? Explique los pasos de la solución.
    \item ¿Cuál es la longitud de onda más larga del espectro de emisión de estos átomos?
    \item ¿En qué parte del espectro electromagnético se encuentra y
    \item a qué serie espectral pertenece?
    \item ¿Cuál es la longitud de onda más corta?
    \item ¿En qué parte del espectro electromagnético se ubica y
    \item a qué serie pertenece?
\end{enumerate}